\documentclass{article}
\usepackage[utf8]{inputenc}
\usepackage[T1]{fontenc}
\usepackage{geometry}
\geometry{
    a4paper,
    left=2.5cm,
    right=2.5cm,
    top=2.5cm,
    bottom=2.5cm
}
\usepackage{graphicx}
\usepackage{float}
\usepackage{hyperref}
\usepackage[english]{babel}
\usepackage{amsmath}
\usepackage{amssymb}
\usepackage{array}
\usepackage{booktabs}
\usepackage{caption}
\usepackage{subcaption}
\usepackage{enumitem}
\usepackage{fancyhdr}
\usepackage{titlesec}
\usepackage{url}
\usepackage{textcomp}
\usepackage{xcolor}
\usepackage{listings}
\usepackage{multirow}
\usepackage{longtable}
\usepackage{tabularx}
\usepackage{times}

% Set up page style
\pagestyle{fancy}
\fancyhf{}
\rfoot{\thepage}
\renewcommand{\headrulewidth}{0pt}

% Title formatting
\titleformat{\chapter}[display]{\normalfont\bfseries}{\chaptertitle\thechapter}{0pt}{\Huge}
\titlespacing*{\chapter}{0pt}{50pt}{40pt}

% Document info
\title{Template Document}
\author{Academic Support Agent}
\date{\today}

\begin{document}
\maketitle

\pretext % n o apague esta linha

% IN CIO DA PARTE TEXTUAL                                           

\section{Introdu o}

Lorem ipsum dolor sit amet, consectetuer adipiscing elit. Ut purus elit, vestibulum ut, placerat ac,adipiscing vitae, felis. Curabitur dictum gravida mauris. Nam arcu libero, nonummy eget, consectetuerid, vulputate a, magna. Donec vehicula augue eu neque. Pellentesque habitant morbi tristique senectuset netus et malesuada fames ac turpis egestas.  Mauris ut leo.  Cras viverra metus rhoncus sem.  Nullaet lectus vestibulum urna fringilla ultrices. Phasellus eu tellus sit amet tortor gravida placerat. Integersapien est, iaculis in, pretium quis, viverra ac, nunc. Praesent eget sem vel leo ultrices bibendum. Aeneanfaucibus. Morbi dolor nulla, malesuada eu, pulvinar at, mollis ac, nulla. Curabitur auctor semper nulla. Donec varius orci eget risus. Duis nibh mi, congue eu, accumsan eleifend, sagittis quis, diam. Duis egetorci sit amet orci dignissim rutrum.

\begin{subequations}
\label{eq:Maxwell}
Equa es de Maxwell:
\begin{align}
        B'&=-\nabla \times E,         \label{eq:MaxB} \\
        E'&=\nabla \times B - 4\pi j, \label{eq:MaxE}
\end{align}
\end{subequations}

Nam dui ligula, fringilla a, euismod sodales, sollicitudin vel, wisi. Morbi auctor lorem non justo\footnote{Exemplo de nota de rodap  }. Nam lacus libero, pretium at, lobortis vitae, ultricies et, tellus.  Donec aliquet, tortor sed accumsanbibendum, erat ligula aliquet magna, vitae ornare odio metus a mi. Morbi ac orci et nisl hendrerit mollis.Suspendisse ut massa. Cras nec ante. Pellentesque a nulla\footnote{Exemplo de mais uma nota de rodap  com um texto bem mais longo de modo que ocupe mais que uma linha Cum sociis natoque penatibus et magnis disparturient montes, nascetur ridiculus mus. Aliquam tincidunt urna.Cum sociis natoque penatibus et magnis disparturient montes, nascetur ridiculus mus. Aliquam tincidunt urna.Cum sociis natoque penatibus et magnis disparturient montes, nascetur ridiculus mus. Aliquam tincidunt urna. Cum sociis natoque penatibus et magnis disparturient montes, nascetur ridiculus mus. Aliquam tincidunt urna.Cum sociis natoque penatibus et magnis disparturient montes, nascetur ridiculus mus. Aliquam tincidunt urna.Cum sociis natoque penatibus et magnis disparturient montes, nascetur ridiculus mus. Aliquam tincidunt urna.}. Cum sociis natoque penatibus et magnis disparturient montes, nascetur ridiculus mus. Aliquam tincidunt urna. Nulla ullamcorper vestibulum turpis. 
\begin{citacao}
 Quisque ullamcorper placerat ipsum. Cras nibh. Morbi vel justo vitae lacus tinciduntultrices.  Lorem ipsum dolor sit amet, consectetuer adipiscing elit.  In hac habitasseplatea dictumst.  Integer tempus convallis augue.  Etiam facilisis.  Nunc elementumfermentum wisi.  Aenean placerat.  Ut imperdiet, enim sed gravida sollicitudin, felisodio placerat quam, ac pulvinar elit purus eget enim. Nunc vitae tortor. Proin tempusnibh sit amet nisl. Vivamus quis tortor vitae risus porta vehicula.
\end{citacao}

\lipsum[5]

\section{Elementos flutuantes}
Tabelas, quadros e figuras s o elementos utilizados em trabalhos acad micos, e possuem diferen as e especifica es definidas pela ABNT.

\subsection{Exemplo de subse o}
As tabelas s o formadas por linhas verticais, devem manter suas bordas laterais abertas e geralmente s o utilizadas para dados quantitativos. Os quadros, por outro lado, s o formados por linhas verticais e horizontais, devem ter todas suas extremidades fechadas e s o mais utilizados para dados qualitativos.

Por  ltimo, as figuras s o elementos ilustrativos, que podem ser em forma de fotos, mapas, gr ficos, gravuras, etc.

\begin{table}[H]
\IBGEtab{
  \caption{Pessoas Residentes em domic lios particulares, por sexo e situa o do domic lio - Brasil - 1980}%
  \label{tab: residentes}
  }
{
\begin{tabular}{m{2.0cm}ccc}
 \toprule
 Situa o do Total & Total & Mulheres & Homens \\
\midrule
Total & 117.960.301 & 59.595.332 & 58.364.969 \\
Urbana & 79.972.931 & 41.115.439 & 38.857.492 \\
Rural & 37.987.370 & 18.479.893 & 19.507.477 \\
\bottomrule
\end{tabular}
}
{
\fonte{Elaborada pelo autor}
\nota{Dados fict cios}
}
\end{table}
		
		
\begin{quadro}[H]
	\centering
	\caption{Alguns detalhes sobre elementos flutuantes}
	\label{quadroex1}
	\begin{tabular}{|l|m{4cm}|m{4cm}|m{4cm}|}
		\hline
		& \textbf{Tabela}                                                                             & \textbf{Quadro}                                                                             & \textbf{Figuras}                                                                    \\ \hline
		\textbf{Formato}    & Bordas laterais n o podem ser fechadas.                                                     & As extremidades devem ser fechadas.                                                         & Podem ser em forma de fotos, mapas, gr ficos, gravuras, etc.                        \\ \hline
		\textbf{Uso}        & Geralmente para dados quantitativos.                                                        & Geralmente para dados qualitativos.                                                         & Ilustrar informa es e dados.                                                       \\ \hline
		\textbf{Elementos}  & T tulo, cabe alho, conte do, fonte e, se necess rio, notas explicativas.                    & T tulo, fonte, legenda e notas.                                                             & T tulo, numera o e fonte.                                                          \\ \hline
		\textbf{Divis o}    & Formada por linhas horizontais.                                                               & Formado por linhas horizontais e verticais.                                                 & -                                                                                   \\ \hline
		\textbf{Formata o} & O n mero e o t tulo da tabela devem vir acima dela, enquanto a fonte deve aparecer embaixo. & O n mero e o t tulo do quadro devem vir acima dele, enquanto a fonte deve aparecer embaixo. & O n mero e o t tulo devem aparecer no topo, enquanto a fonte deve aparecer embaixo. \\ \hline
	\end{tabular}
	\fonte{Elaborado pelo autor}
\end{quadro}
				
Abaixo, a figura \ref{fig:ex}   apresentada como exemplo de disposi o dos elementos legenda, numera o e fonte.
Para uso avan ado de imagens no \LaTeX, recomenda-se a consulta de literatura especializada \cite{Goossens2007}.

	\begin{figure}[H]
		\centering
		\caption{Exemplo de inclus o de imagem}
		\includegraphics[height=6cm]{./figuras/wave.png}
		\fonte{\textcite{opensource}}
		\label{fig:ex}
	\end{figure}

\subsection{Outros exemplos}
Nesta subse o apresentamos alguns elementos deste modelo em \LaTeX.

\subsubsection{Cita es diretas com mais de tr s linhas}
Neste modelo, para fazermos uma cita o direta com mais de tr s linhas, podemos usar o ambiente \texttt{citacao}. Por exemplo, se digitarmos
\begin{verbatim}
\begin{citacao}
    Lorem ipsum dolor sit amet, consectetur adipiscing elit.
    Nulla in pulvinar nulla. Sed tempor condimentum felis quis blandit.
    Maecenas ultricies sapien eu nisi hendrerit maximus.
    Nullam sagittis consectetur leo, at feugiat lorem.
    Nulla sed maximus metus, sed tincidunt leo.
    Fusce fermentum consectetur libero, eget rhoncus nisl laoreet eget.
    Nam eget justo metus.
\end{citacao}
\end{verbatim}
obteremos o seguinte resultado
\begin{citacao}
    Lorem ipsum dolor sit amet, consectetur adipiscing elit. Nulla in pulvinar nulla. Sed tempor condimentum felis quis blandit. Maecenas ultricies sapien eu nisi hendrerit maximus. Nullam sagittis consectetur leo, at feugiat lorem. Nulla sed maximus metus, sed tincidunt leo. Fusce fermentum consectetur libero, eget rhoncus nisl laoreet eget. Nam eget justo metus.
\end{citacao}

Lorem ipsum dolor sit amet, consectetur adipiscing elit. Nulla in pulvinar nulla. Sed tempor condimentum felis quis blandit.

%--------------------------------
% Parte p s-textual
%--------------------------------

\section*{Agradecimentos}
Esta   uma se o reservada aos agradecimentos dos autores, caso pertinente.	
\cite{Lei:12651:2012}
% REFER NCIAS BIBLIOGR FICAS

\singlespacing
\printbibliography
\end{document}